\subsection{UPPER LIMIT AND LOWER LIMIT OF A REAL-VALUED FUNCTION WITH RESPECT TO A FILTER}

Let \(f\) be a real-valued function defined on a set \(\mathbf{X}\) filtered by a filter \(\mathfrak{G}\) . \(\mathfrak{G}\) is a directed set with respect to the relation \(\supset\) (Chapter I, § 6). For each \(\mathbf{M} \in \mathfrak{G}\) consider the real number \(\sup_{x \in \mathbf{M}} f(x)\) : we have a function \(\mathbf{M} \to \sup_{x \in \mathbf{M}} f(x)\) of \(\mathfrak{G}\) into \(\overline{\mathbf{R}}\) , which is a decreasing function on \(\mathfrak{G}\) , by Proposition 7 of no. 4. Hence, by Theorem 2 of no. 2, it has a limit with respect to the directed set \(\mathfrak{G}\) .

DEFINITION 6. The limit of the real-valued function \(\mathbf{M} \to \sup_{x \in \mathbf{M}} f(x)\) with respect to the directed set \(\mathfrak{G}\) is called the upper limit of \(f\) with respect to the filter \(\mathfrak{G}\) , and is denoted by \(\lim \sup_{\mathfrak{G}} f\) , or by \(\lim \sup_{x, \mathfrak{G}} f(x)\) .

The lower limit of \(f\) with respect to the filter \(\mathfrak{G}\) is defined similarly, and is denoted by \(\liminf_{\mathfrak{G}} f\) or \(\liminf_{x, \mathfrak{G}} f(x)\) . Thus we have

\[
\left\{ \begin{array}{l} \lim \sup_{\mathfrak{G}} f = \lim_{\mathbf {M} \in \mathfrak {G}}(\sup_{x \in \mathbf {M}}f(x)), \\ \lim \inf_{\mathfrak{G}} f = \lim_{\mathbf {M} \in \mathfrak {G}}(\inf_{x \in \mathbf {M}}f(x)). \end{array} \right.\tag{10}
\]

Often the filter \(\mathfrak{G}\) is suppressed from the notation, and we write simply \(\lim \sup f\) or \(\lim \sup_x f(x)\) , or \(\lim \sup f(x)\) when there is no risk of confusion.

From formulae (10) and Theorem 1 we have

\[
\inf _ {x \in \mathbf {X}} f (x) \leqslant \lim \inf _ {\mathfrak {G}} f \leqslant \lim \sup _ {\mathfrak {G}} f \leqslant \sup _ {x \in \mathbf {X}} f (x).\tag{11}
\]

By Theorem 2 of no. 2 we may also write

\[
\left\{\lim \sup _ {\mathfrak {G}} f = \inf _ {\mathbf {M} \in \mathfrak {G}} (\sup _ {x \in \mathbf {M}} f (x)), \right.\tag{12}
\]

\[
\left\{\lim \inf _ {\mathfrak {G}} f = \sup _ {\mathbf {M} \in \mathfrak {G}} (\inf _ {x \in \mathbf {M}} f (x)). \right.
\]

Also we may replace the filter \(\mathfrak{G}\) , on the right-hand sides of formulae (10) and (12), by any base \(\mathfrak{B}\) of \(\mathfrak{G}\) .

From (2) and (10),

\[
\lim \inf _ {\mathfrak {G}} f = - \lim \sup _ {\mathfrak {G}} (- f)\tag{13}
\]

and therefore we need consider only the upper limit.

THEOREM 3. The upper limit of a real-valued function \(f\) with respect to a filter \(\mathfrak{G}\) is equal to the largest cluster value of \(f\) with respect to \(\mathfrak{G}\) .

Let \(b\) be a cluster point of \(f\) with respect to \(\mathfrak{G}\) . For each \(\mathbf{M} \in \mathfrak{G}\) , \(b\) lies in the closure of \(f(\mathbf{M})\) , hence \(b \leqslant \sup_{x \in \mathbf{M}} f(x)\) , and therefore, by (12), \(b \leqslant \limsup_{\mathfrak{G}} f = a\) .

On the other hand, let V be any open neighbourhood of a in \(\overline{R}\) . Then there exists a set \(M_{0}\) in G such that, for each \(M \in G\) contained in \(M_{0}\) , we have \(\sup_{x \in M} f(x) \in V\) ; since V is open it follows that \(f(M)\) meets V, and therefore a is a cluster point of f with respect to G, and the proof is complete.

COROLLARY 1. In order that \(\limsup_{\mathfrak{G}}f = \liminf_{\mathfrak{G}}f\) , it is necessary and sufficient that \(f\) has a limit with respect to the filter \(\mathfrak{G}\) , and then

\[
\lim _ {\mathfrak {G}} f = \lim \sup _ {\mathfrak {G}} f = \lim \inf _ {\mathfrak {G}} f.
\]

For since \(\overline{\mathbf{R}}\) is compact, the filter base \(f(\mathfrak{G})\) has a limit point if and only if it has only one cluster point (Chapter I, § 9, no. 1, Corollary to Theorem 1).

COROLLARY 2. If \(\mathfrak{H}\) is a filter finer than \(\mathfrak{G}\) , we have

\[
\lim \inf _ {\mathfrak {G}} f \leqslant \lim \inf _ {\mathfrak {H}} f \leqslant \lim \sup _ {\mathfrak {H}} f \leqslant \lim \sup _ {\mathfrak {G}} f.
\]

For every cluster point of \(f\) with respect to \(\mathfrak{H}\) is also a cluster point of \(f\) with respect to \(\mathfrak{G}\) (Chapter I, § 7, no. 3).

In particular, if \(\lim_{\mathfrak{H}}f\) exists, then

\[
\lim \inf _ {\mathfrak {G}} f \leqslant \lim _ {\mathfrak {H}} f \leqslant \lim \sup _ {\mathfrak {G}} f.
\]

COROLLARY 3. Let A be a set of the filter G, let \(\mathfrak{G}_{\mathbf{A}}\) be the filter induced on A by G, and let \(f_{\mathbf{A}}\) be the restriction of \(f\) to A; then

\[
\lim \sup _ {\mathfrak {G} _ {\mathbf {A}}} f _ {\mathbf {A}} = \lim \sup _ {\mathfrak {G}} f.
\]

For every cluster point of the filter base \(f(\mathfrak{G})\) is a cluster point of the filter base \(f_{\mathrm{A}}(\mathfrak{G}_{\mathrm{A}})\) , and conversely.

For this reason, if \(f\) is defined only on a subset \(A\) of \(X\) belonging to \(\mathfrak{G}\) , we shall often write \(\limsup_{\mathfrak{G}} f\) instead of \(\limsup_{\mathfrak{G}_{\mathbf{A}}} f_{\mathbf{A}}\) , by abuse of language.

PROPOSITION 11. Let \(f\) and \(g\) be two real-valued functions defined on a filtered set \(X\) . Then the relation \(f \leqslant g\) implies

\[
\left\{ \begin{array}{l} \lim \sup f \leqslant \lim \sup g, \\ \lim \inf f \leqslant \lim \inf g. \end{array} \right.\tag{14}
\]

This is an immediate consequence of the relations (12).

When X is a topological space and \(\mathfrak{G}\) is the neighbourhood filter of a point \(a\) of X, we write \(\limsup_{x\to a}f(x)\) {[}resp. \(\liminf_{x\to a}f(x)]\) in place of \(\limsup_{\mathfrak{G}}f\) {[}resp. \(\liminf_{\mathfrak{G}}f]\) ; clearly we have

\[
\liminf _ {x \to a} f (x) \leqslant f (a) \leqslant \limsup _ {x \to a} f (x).\tag{15}
\]

More generally, if X is a subspace of a topological space Y, and if G is the trace on X of the neighbourhood filter of a point \(a \in \overline{X}\) , we write \(\limsup_{x \to a, x \in X} f(x)\) {[}resp. \(\liminf_{x \to a, x \in X} f(x)\) {]} instead of \(\limsup_{\mathfrak{G}} f\) {[}resp. \(\liminf_{\mathfrak{G}} f\) {]}; \(\limsup_{x \to a, x \in X} f(x)\) is called the upper limit of \(f(x)\) as x tends to a while remaining in X. If X is the complement of \(\{a\}\) we write `` \(x \neq a\) '' in place of `` \(x \in X\) '' in these notations.

If \(\mathbf{A}\) is a subset of \(\mathbf{X}\) such that \(a \in \overline{\mathbf{A}}\) , then (Corollary 2 to Theorem 3)

\[
\liminf _ {x \to a, x \in \mathbf {X}} f (x) \leqslant \liminf _ {x \to a, x \in \mathbf {A}} f (x) \leqslant \limsup _ {x \to a, x \in \mathbf {A}} f (x) \leqslant \limsup _ {x \to a, x \in \mathbf {X}} f (x).
\]

If V is a neighbourhood of \(a\) in Y, we have (Corollary 3 to Theorem 3)

\[
\lim \sup _ {x \to a, x \in \mathbf {V} \cap \mathbf {X}} f (x) = \lim \sup _ {x \to a, x \in \mathbf {X}} f (x).
\]

Hence the notions of upper and lower limits at a point of a topological space are, like the notion of limit, of local character.

Finally, if \(\mathfrak{G}\) is the Fréchet filter on \(\mathbf{N}\) , the upper (resp. lower) limit, with respect to \(\mathfrak{G}\) , of the mapping \(n \to u_n\) of \(\mathbf{N}\) into \(\overline{\mathbf{R}}\) is denoted by \(\limsup_{n \to \infty} u_n\) (resp. \(\liminf_{n \to \infty} u_n\) ) and is called the upper (resp. lower) limit of the sequence of real numbers \(u_n\) .

The relation \(\limsup_{n\to \infty}u_n = a\in \mathbb{R}\) is therefore equivalent to the following: given any \(\varepsilon >0\) there exists an integer \(n_0\) such that, for each \(n\geqslant n_0\) we have \(u_{n}\leqslant a + \varepsilon\) , and for an infinity of values of \(n\) we have \(u_{n}\geqslant a - \varepsilon\) . The definition of the upper limit of a sequence may be translated similarly when its value is \(+\infty\) or \(-\infty\) .

Given a sequence \((f_n)\) of real-valued functions defined on a set \(\mathbf{X}\) , we denote by \(\limsup_{n\to \infty}f_n\) (resp. \(\liminf_{n\to \infty} f_n\) ) the real-valued function whose value at any point \(x\in \mathbf{X}\) is \(\limsup_{n\to \infty}f_n(x)\) {[}resp. \(\liminf_{n\to \infty} f_n(x)\) {]}. From (10) and (12) we deduce

\[
\left\{\begin{array}{l}\lim \sup _ {n \rightarrow \infty} f _ {n} = \inf _ {n \in \mathbb {N}} (\sup _ {m \geqslant n} f _ {m}) = \lim _ {n \rightarrow \infty} (\sup _ {m \geqslant n} f _ {m}),\\\lim \inf _ {n \rightarrow \infty} f _ {n} = \sup _ {n \in \mathbb {N}} (\inf _ {m \geqslant n} f _ {m}) = \lim _ {n \rightarrow \infty} (\inf _ {m \geqslant n} f _ {m}),\end{array}\right.\tag{16}
\]

the limits being taken in the product space \(\overline{\mathbf{R}}^{\mathbf{X}}\) . The sequence \((f_n)\) has a limit in \(\overline{\mathbf{R}}^{\mathbf{X}}\) if and only if \(\limsup_{n\to\infty}f_n=\liminf_{n\to\infty}f_n\) (Corollary 1 to Theorem 3, and Chapter I, § 7, no. 6, Corollary 1 to Proposition 10).

